% !TeX root = ../../Book3.tex
\section{Picard 大定理}
\label{b3:sec:1405}

本性奇点附近的稠密像由 Casorati--Weierstrass 给出，Picard 大定理进一步保证除至多一值外，每个值在任意穿孔邻域取到无穷次。正规族把局部遗漏转成整平面极限，是增强结论的桥梁。

\subsection{本性奇点}
\label{b3:subsec:140501}

设 \(f\) 在 \(0<|z-a|<R\) 全纯且 \(a\) 为本性奇点。这里本性是 Laurent 主部无限，不是简单的无界性；极点同样无界，却在小邻域遗漏一个大圆盘中的全部值。

\ProseParagraphBreak
若某个值在一个穿孔邻域只有有限个原像，缩小邻域后便可完全遗漏该值。因此证明大定理只需证明本性奇点附近不能同时遗漏两个不同有限值，再把这个结论用于有限原像的两个候选例外。

\subsection{Picard 大定理}
\label{b3:subsec:140502}

\term[picard-da-dingli]{Picard 大定理}[great Picard theorem]限制本性奇点附近的例外值。

\begin{theorem}[Picard 大定理]\label{b3:thm:14-picard-great}
在全纯函数的本性孤立奇点附近，除至多一个有限值外，每个复值在任意穿孔邻域中取得无穷次。
\end{theorem}
\begin{proof}
平移奇点到0。假设某穿孔盘上遗漏两个值，仿射归一化为0与1。选 \(b\notin\{0,1\}\)，Casorati--Weierstrass 给 \(z_n\rightarrow0\) 且 \(f(z_n)\rightarrow b\)。重缩放 \(g_n(\zeta)=f(z_n\zeta)\) 在 \(0<|\zeta|<R/|z_n|\) 定义，遗漏0、1、无穷远。对 \(\mathbb C^*\) 的紧穷竭用省略值 Montel 和对角抽取，得到球面极限 \(g\)，且 \(g(1)=b\)。

若 \(g\) 非恒常，它仍遗漏三值；复合指数 \(g(e^w)\) 是全平面非恒常亚纯函数且遗漏三值，与 Picard 小定理矛盾。故 \(g\equiv b\)。由于这个常数有限，收敛在单位圆上是普通一致，充分晚有 \(|f(z_n\zeta)|\le M\)、\(|\zeta|=1\)。抽取 \(|z_n|\) 严格递减子列，这给趋于0的一列完整圆周上的同一界。相邻圆周之间 \(f\) 全纯且连续到边界，最大模原理使整个闭圆环有同一界。这些圆环覆盖一个穿孔邻域，所以 \(f\) 有界，0可去，矛盾。

因此两个遗漏值不可能。若两个不同值各只有有限次原像，缩小穿孔盘同时避开这些有限原像，便得到两个遗漏值，同样矛盾，故至多一值例外。
\end{proof}
对亚纯本性奇点，至多两个球面值例外；若有三个，先把其中一个送到无穷远，使小穿孔盘上函数全纯，再用刚证的版本。

\subsection{与 Casorati–Weierstrass 定理的比较}
\label{b3:subsec:140503}

\(e^{1/z}\) 在0本性且遗漏0；对每个 \(b\ne0\)，解为 \(z=(\log b+2\pi\mathrm{i}n)^{-1}\)，除去可能分母为0的一项，得到无穷原像趋于0。这个例子使全部量词都可直接核对：每个非零值、每个穿孔邻域、无穷多个原像。

\ProseParagraphBreak
Casorati--Weierstrass 只保证像稠密，从接近某值不能推出真的取到该值。大定理的证明通过正规族、覆盖与最大模原理填补这个差距，不把“稠密”误写成“满射”。大定理也不声称例外值必存在，某些本性函数没有有限例外。

\ProseParagraphBreak
大 Picard 中的“无限次”可从缩小邻域的形式严格推出。若某值 \(a\) 在一个穿孔邻域只有有限个原像，就把圆盘缩到排除全部这些原像；它便在更小穿孔邻域省略 \(a\)。若有两个不同值各只取有限次，缩到同时排除两组原像便得到两个省略值，违反全纯大 Picard。亚纯版本同理，但容许两个球面例外值而非一个有限例外值。

\ProseParagraphBreak
Casorati--Weierstrass 的稠密像不保证每个指定值真的出现。一个开稠密集仍可以缺少无限多个点；本性奇点的稠密性只说目标任何小邻域都有像。Picard 则把未取值的集合压缩到至多一个或两个点。证明中利用稠密性选择一个通常值的逼近序列，再利用正规族把局部逼近转成一整圈的统一控制，最后用最大模原理推出可去性。稠密性、紧致性和环带估计各自承担不同步骤，缺少其中之一就不能从稠密像升级到省略值结论。
